\chapter{The product: scenes, projects, actions and the vendable}\label{ch:product}

\begin{onething}
Allotrope is a web application with four nouns (Scene, Project, Action, Export). Every piece of
analysis in this book runs as an Action, Actions link into a directed acyclic graph, and every
Action consumes one standard object, the \emph{vendable}: a cube in physical units, a validity
cube of the same shape, and per-band bookkeeping, with no map coordinates at all.
\end{onething}

\section{Why this comes before any model}\label{sec:product:motivation}

You are about to spend many chapters on detectors and neural networks. If you have only ever seen
machine-learning code, you will picture a script that takes a file and prints a result, and you
will then be unable to work out why the code is shaped the way it is: why a model is ``one option
inside a step'', why a file called \code{filtered_vendable.pkl} must exist before anything can be
scored, why renaming a class can break every stored result. Each of those follows from the
product's structure, so this chapter lays the structure out first.

A user logs in through a browser, uploads a satellite file, clicks through a chain of steps,
drags a threshold slider, and downloads a result. Behind that sits a Postgres database, a login
system, a React front end and a background worker process. The models are components inside it.

\section{The four nouns}\label{sec:product:nouns}

\begin{definition}[Scene, Project, Action, Export]\index{Scene}\index{Project}\index{Action}\index{Export}
A \textbf{Scene} is one uploaded raw satellite acquisition, onboarded into the standard form of
\cref{sec:product:vendable}. A \textbf{Project} is a piece of work bound to exactly one Scene. An
\textbf{Action} is one processing step inside a Project; it reads the outputs of earlier Actions
and writes its own outputs to its own folder (an \textbf{ActionOutput}). An \textbf{Export} is
the deliverable: a zip bundle with a GeoTIFF, a Shapefile and a CSV naming every candidate anomaly
with its coordinates.
\end{definition}

The relationships are worth stating exactly, because they are asked about constantly.
\begin{itemize}
\item One Scene can back many Projects; a Project binds to exactly one Scene. Analysing one scene
two ways means one Scene and two Projects. Comparing two scenes means two Projects.
\item A Project holds many Actions. An Action names its inputs by referring to earlier Actions'
outputs.
\end{itemize}

On the wire (in every API request and response), identifiers are never bare database keys. They
carry a type prefix, \code{scene_<uuid>}, \code{action_<uuid>}, \code{output_<uuid>}, where a
\textbf{UUID}\index{UUID} is a $128$-bit random identifier written as $32$ hexadecimal digits.
The helpers \code{to_prefixed} and \code{parse_prefixed_id} are in
\file{backend/allotrope/api/wireformat.py}. A prefixed id cannot be passed where an id of another
type is expected without the mistake being caught.

\section{Actions form a directed acyclic graph}\label{sec:product:dag}

\subsection{Graphs, briefly}

\begin{definition}[Directed graph, path, cycle, DAG]\index{directed acyclic graph}\index{graph}
A \textbf{directed graph} is a set of \textbf{nodes} and a set of \textbf{edges}, each edge an
arrow from one node to another. A \textbf{path} follows arrows from node to node. A
\textbf{cycle} is a path that returns to where it started. A \textbf{directed acyclic graph}
(DAG) is a directed graph with no cycles.
\end{definition}

In a Project, each Action is a node, and there is an edge from Action $P$ to Action $Q$ when $Q$
reads an output of $P$. The graph is acyclic by construction: an Action can only name outputs that
already exist when it is created, so an arrow always points from an older Action to a newer one,
and a path that only moves forward in time can never come back to its start.

\subsection{The shipped chain, read from the source}

Each Action type declares the inputs it accepts in its \code{META} object, as a list of
\code{ActionInputSpec}s whose \code{producing_action_types} field says which earlier types may
feed it (\file{backend/allotrope/action_types/_meta.py}). Reading those declarations gives the
edges of \cref{fig:product:dag}:
\begin{itemize}
\item \code{band_filter_apply} reads the Scene itself and writes \code{filtered_vendable.pkl}.
\item \code{scene_segmentation} (hyperspectral) reads a \code{band_filter_apply} output and
writes \code{keep_mask.tif}, the map of pixels to keep (not cloud, not shadow).
\code{cloud_mask} plays the same role for thermal scenes.
\item \code{anomaly_scoring} reads a \code{band_filter_apply} output (required for hyperspectral
scenes) and, optionally, a \code{scene_segmentation} or \code{cloud_mask} output and a
ground-truth annotation (\file{anomaly_scoring.py}, lines 250--296). It writes one folder per
model, \code{models/<codename>/}, holding \code{anomaly_score.tif} and \code{reconstruction.tif}
among others.
\item \code{anomaly_detection_prep} reads an \code{anomaly_scoring} output, rescales and combines
the per-model score maps into \code{composite_score.tif}, and stops in the state
\code{needs_threshold} (below).
\item \code{spectral_library_match} (hyperspectral only) reads a committed
\code{anomaly_detection_prep} output and names the material of each anomalous pixel.
\item Export (\code{POST /actions/\{id\}/export}) bundles the result, recovering map coordinates
on the way (\cref{sec:product:georef}).
\end{itemize}

\begin{figure}[htbp]\centering
\begin{tikzpicture}[
  act/.style={draw,rounded corners=2pt,align=center,font=\footnotesize,minimum width=3.3cm,minimum height=0.95cm,fill=white},
  outl/.style={font=\scriptsize\ttfamily,text=spec!80!black},
  arr/.style={-{Stealth[length=2mm]},thick}]
\node[act,draw=spec,fill=spec!8] (scene) at (0,0) {Scene\\[-1pt]{\scriptsize raw file + onboarding vendable}};
\node[act,draw=spec] (bf) at (0,-1.7) {\texttt{band\_filter\_apply}};
\node[outl,right=2pt of bf] {filtered\_vendable.pkl};
\node[act] (seg) at (-3.2,-3.4) {\texttt{scene\_segmentation}\\[-1pt]{\scriptsize or \texttt{cloud\_mask} (thermal)}};
\node[outl,below=1pt of seg] {keep\_mask.tif};
\node[act,draw=spat,fill=spat!8] (sc) at (0,-5.1) {\texttt{anomaly\_scoring}\\[-1pt]{\scriptsize models ticked: Indradhanu, MNF-RX, \dots}};
\node[outl,right=2pt of sc,align=left] {models/<codename>/\\anomaly\_score.tif};
\node[act,draw=loss,fill=loss!8] (prep) at (0,-6.9) {\texttt{anomaly\_detection\_prep}\\[-1pt]{\scriptsize waits in \texttt{needs\_threshold}}};
\node[outl,right=2pt of prep,align=left] {composite\_score.tif\\anomaly\_mask.tif (on commit)};
\node[act] (slm) at (0,-8.7) {\texttt{spectral\_library\_match}\\[-1pt]{\scriptsize hyperspectral only}};
\node[outl,right=2pt of slm,align=left] {family\_matches.parquet, family\_map.tif\\matches.parquet, match\_map.tif};
\node[act,draw=maskgrey] (ex) at (0,-10.3) {Export\\[-1pt]{\scriptsize zip: GeoTIFF + Shapefile + CSV}};
\node[act,draw=maskgrey,dashed] (gt) at (-4.1,-6.9) {annotation\\[-1pt]{\scriptsize ground truth (optional)}};
\draw[arr] (scene) -- (bf);
\draw[arr] (bf) -| (seg);
\draw[arr] (bf) -- (sc);
\draw[arr] (seg) |- (sc);
\draw[arr,dashed,maskgrey] (gt) -- (sc);
\draw[arr] (sc) -- (prep);
\draw[arr] (prep) -- (slm);
\draw[arr] (slm) -- (ex);
\draw[arr,maskgrey] (scene.west) -- (-6.3,0) -- node[pos=0.5,sloped,above,font=\scriptsize]{raw file re-read for coordinates at export} (-6.3,-10.3) -- (ex.west);
\end{tikzpicture}
\caption{The shipped Action graph, edges read from each Action type's declared inputs. Each box
writes its outputs (monospace) to its own folder. \texttt{anomaly\_scoring} has two possible
parents besides the optional annotation, so the structure is a graph, not a chain; Indradhanu is
one model ticked \emph{inside} \texttt{anomaly\_scoring}. Export re-reads the raw scene file to
recover map coordinates.}\label{fig:product:dag}
\end{figure}

\coursefig[\linewidth]{iw-05-chain}{The same chain as drawn in the Indradhanu course (part 5),
emphasising that Indradhanu is one option inside the \texttt{anomaly\_scoring} Action, not an
Action of its own.}{fig:product:coursechain}

\begin{sourcenote}[: the chain's outputs in the source]
The course figure lists \code{matches.parquet} as the output of \code{spectral_library_match}.
The current source writes that file and \code{match_map.tif}, and also the per-family answer
\code{family_matches.parquet} and \code{family_map.tif}
(\file{backend/allotrope/action_types/_spectral_library_match_run.py}, lines 46--47); the
project instructions call the family map ``the answer''. The companion volume on spectral matching
covers both. Second, the docstring of \file{anomaly_detection_prep.py} (line 22) says the
threshold is locked by ``a separate \code{anomaly_detection_commit} action''. In the current
source the commit is an API endpoint, \code{POST /actions/\{id\}/anomaly_detection_commit}, in
\file{backend/allotrope/api/action_threshold.py}; it writes \code{anomaly_mask.tif}, which
\code{spectral_library_match} then reads.
\end{sourcenote}

\subsection{An order in which a graph can run: topological sorting}

A DAG always has at least one order in which every node comes after all of its parents, called a
\textbf{topological order}\index{topological order}. Such an order is what lets a user (or a
script) create Actions one at a time, each with its inputs already finished. The standard way to
find one is Kahn's algorithm: repeatedly pick a node with no unfinished parents, list it, and
remove its outgoing arrows.

\begin{worked}[label={we:product:topo}]{A topological order for one Project}
A Project has six Actions: $B$ (\code{band_filter_apply}), $G$ (\code{scene_segmentation}),
$S$ (\code{anomaly_scoring}), $P_1$ and $P_2$ (two \code{anomaly_detection_prep} Actions trying
different weightings on the same scoring run) and $M$ (\code{spectral_library_match}, from
$P_2$). The edges are
\[ B \to G,\quad B \to S,\quad G \to S,\quad S \to P_1,\quad S \to P_2,\quad P_2 \to M. \]
Count each node's parents (its \textbf{in-degree}):
\[ B{:}\,0,\quad G{:}\,1,\quad S{:}\,2,\quad P_1{:}\,1,\quad P_2{:}\,1,\quad M{:}\,1. \]
Run Kahn's algorithm on paper.
\begin{enumerate}
\item Only $B$ has in-degree $0$. List $B$; remove $B \to G$ and $B \to S$. Now
$G{:}\,0$, $S{:}\,1$.
\item Only $G$ has in-degree $0$. List $G$; remove $G \to S$. Now $S{:}\,0$.
\item List $S$; remove $S \to P_1$ and $S \to P_2$. Now $P_1{:}\,0$, $P_2{:}\,0$.
\item Two choices. List $P_1$ (it has no outgoing arrows).
\item List $P_2$; remove $P_2 \to M$. Now $M{:}\,0$.
\item List $M$.
\end{enumerate}
One valid order is $B, G, S, P_1, P_2, M$. Another is $B, G, S, P_2, M, P_1$: step 4 had a
choice, so topological orders are not unique. What every order shares is that $S$ comes after
both $B$ and $G$, because $S$ has two parents.
\end{worked}

There is no scheduler and no automatic re-run in the product: a user creates each Action. The graph
is the record of what was derived from what, which is why scoring once and trying several
thresholds (two prep Actions sharing one scoring parent, as above) costs one expensive run.

\section{Two processes and one queue}\label{sec:product:queue}

Two programs run. The \textbf{api} is the web server: it handles requests, reads and writes the
database, and hands out work, but never does heavy computation. The \textbf{worker} does the
analysis and has no web interface. They communicate through a \code{jobs} table in Postgres:
there is no Redis, no Celery and no message broker.

\begin{itemize}
\item The api inserts a job row with status \code{queued}.
\item A worker \emph{claims} one with \code{SELECT \dots\ FOR UPDATE SKIP LOCKED}
(\file{backend/allotrope_worker/claim.py}), sets it to \code{running} and commits immediately.
``Skip locked'' means two workers asking at once never get the same row: each skips rows another
transaction is holding.
\item While it works, the worker writes a \textbf{heartbeat}\index{heartbeat} timestamp to the row
at a fixed interval.
\item A \textbf{reaper}\index{reaper} flips any \code{running} job whose last heartbeat is older
than a timeout to \code{failed} (\file{backend/allotrope_worker/reaper.py}), so a crashed worker
does not leave a job ``running'' forever.
\end{itemize}

\begin{worked}[label={we:product:heartbeat}]{When is a job reaped?}
The defaults are a heartbeat every $5$ seconds and a timeout of $60$ seconds\tagM{backend/allotrope_worker/runner.py},
overridable by the environment variables \code{WORKER_HEARTBEAT_INTERVAL_SEC} and
\code{WORKER_HEARTBEAT_TIMEOUT_SEC}. A job's last heartbeat was written at 10:00:00 and the worker
then crashed. The reaper computes a cutoff of ``now minus $60$ s'' and reaps rows whose heartbeat is
older than the cutoff.
\begin{itemize}
\item At 10:00:59 the cutoff is 09:59:59. The heartbeat (10:00:00) is newer: kept.
\item At 10:01:01 the cutoff is 10:00:01. The heartbeat is older: reaped to \code{failed}.
\end{itemize}
How many heartbeats must a healthy worker miss before that happens?
\[ 60 / 5 = 12 \text{ intervals.} \]
The source comment asks for the timeout to be at least about three heartbeat intervals ``to absorb
jitter''; the defaults give twelve.
\end{worked}

The consequence you will feel: \textbf{the api can be perfectly healthy while the worker is dead}.
The api's health check talks to the database, not the worker. This has happened twice in this
repository's history, and jobs sat in \code{queued} forever. When something does not run, check
the worker first.

\subsection{The step that waits for a person}

\code{anomaly_detection_prep} is the one Action type that does not finish in \code{complete}. It
writes its composite score map and stops in \code{needs_threshold}, waiting for a person to drag a
threshold slider in the viewer. Each ``Apply'' recomputes a preview mask; the commit endpoint locks
the threshold and writes \code{anomaly_mask.tif}. \Cref{sec:problem:threshold} gave the reason:
no fixed cut-off works across scenes, so the cut is a human judgement made per scene.

\subsection{Where the code lives, and the lazy-import rule}

\begin{table}[htbp]\centering
\caption{The layers. The rule in the last column is enforced by review, not by tooling.}\label{tab:product:layers}
\small
\begin{tabular}{@{}llp{0.48\linewidth}@{}}\toprule
Layer & Path & Rule\\\midrule
Portable science & \texttt{app/} & numpy, torch, rasterio only. No database, no web framework.\\
Shared backend & \texttt{backend/allotrope/} & database models, routes, authentication, the Action registry; used by api and worker.\\
Worker only & \texttt{backend/allotrope\_worker/} & the claim loop, heartbeat, reaper and job handlers.\\\bottomrule
\end{tabular}
\end{table}

The models' architectures, trainers and inferencers live in \code{app/}. They know nothing about
users, jobs, HTTP or Postgres, which is why the same model code runs under a training script on a
rented GPU machine and under the worker in a container.

Action modules come in pairs: \code{anomaly_scoring.py} (light: name, metadata, configuration
validation) and \code{_anomaly_scoring_run.py} (heavy: the implementation). Action modules must not
import \code{app.*}, \code{torch} or \code{rasterio} at the top of the file, only inside their
functions. The reason: the api imports every Action module at start-up to list the available
Actions, and a top-level \code{import torch} would drag hundreds of megabytes of deep-learning
library into the web server to render a drop-down menu. A side effect worth remembering: a broken
import in a worker-only module does not stop the api starting, so the api stays green while the
worker cannot start.

\section{Where Indradhanu lives}\label{sec:product:models}

Indradhanu is not an Action. It is one selectable model \emph{inside} \code{anomaly_scoring}: the
user ticks one or more models by name in the scoring dialog, and the Action loops over them
(\code{for codename in cfg["model_codenames"]:}), writing each model's outputs to its own
\code{models/<codename>/} folder. You can run Indradhanu and MNF-RX in one Action and compare their
score maps side by side.

Every foundation model carries an Indic codename, and the backend selects models \emph{by
codename}, not by class name.

\begin{table}[htbp]\centering
\caption{The foundation-model roster (Indradhanu course, part 5; project instructions).}\label{tab:product:codenames}
\small
\begin{tabular}{@{}llll@{}}\toprule
Codename & Meaning & Architecture & Sensor\\\midrule
Pratibimba & reflection & spatial autoencoder & thermal\\
Antardhana & disappearance & masked autoencoder & thermal\\
Tirohita & vanished & masked autoencoder, L1 loss & thermal\\
Asanskrita & unrefined & unnormalised variant & thermal\\
Drashta & seer & normalised masked autoencoder & thermal\\
Chakshu & the eye & \texttt{segformer\_mae} & thermal\\
Indradhanu & rainbow (``Indra's bow'') & \texttt{hyperspectral\_segformer\_mae} & hyperspectral\\\bottomrule
\end{tabular}
\end{table}

A codename becomes a running model through three cooperating pieces: a \textbf{manifest}
(\code{current.json}) beside the weights file under \code{allotrope_models/<architecture>/}
saying which file is current; the \textbf{resolver}
(\file{backend/allotrope/foundation_models/resolver.py}), a static table of what each architecture
can do (which scoring methods to offer, default patch size and stride, where its normalisation
statistics live); and the \textbf{inferencer factory}
(\file{app/foundation_models/inferencers/inferencer_factory.py}) mapping the architecture to the
Python class that runs it. Adding a model means four registrations: the architecture class, the
trainer factory, the inferencer factory, and a resolver capabilities entry plus a manifest. The
fourth is the one people forget, and it fails silently: the model catalogue skips architectures
with no capabilities entry, so a model that trained perfectly never appears in the picker.

\begin{drift}
\emph{Manifest values quoted by the Indradhanu course (part 5).} \code{current.json} for
\code{hyperspectral_segformer_mae} listed file
\code{hyperspectral_segformer_mae_v0.2.0_epoch200.pt}, version v0.2.0, epoch $200$, $5{,}507{,}354$
parameters, encoder widths $[32, 64, 160, 256]$, compression width $D = 32$, decoder width $256$,
validation loss $0.04349$\tagM{indradhanu workings/05-where-it-lives.md}. The manifest itself is
not in this checkout (only the \code{.pt} file is), so these could not be re-read for this book.
Since 2026-10-29 the shipped hyperspectral model is Indradhanu V2 (``Indradhanu Heavy''), whose
resolver entry reads native bands (\cref{ch:sensors}, \cref{ch:shipping,ch:heavy}).
\end{drift}

\section{The vendable}\label{sec:product:vendable}

Raw satellite files are unpleasant (\cref{ch:sensors}): every vendor its own layout, quality flags
packed into bits, some files reflectance and some raw counts, detector quirks everywhere. If every
detector had to cope with all of that, every detector would change whenever a sensor was added.
So the first thing the system does with a scene is convert it, once, into a standard object.

\begin{definition}[Vendable]\index{vendable}
A \textbf{vendable} is the standard in-memory form of a scene: a cube in physical units, a
validity cube of the same shape, and per-band metadata. It is produced by a
\textbf{DatasetBuilder}\index{DatasetBuilder} from a raw file. Every detector, model and Action
consumes a vendable; nothing downstream reads a raw sensor file. The classes are in
\file{app/models/dataset/vendables.py}.
\end{definition}

Read the odd name as ``the thing this scene can \emph{vend}'': what it dispenses to anything that
asks, once cleaned up.

\begin{analogy}
A shipping port receives cargo in every kind of packaging (sacks, crates, barrels, loose bales) and
repacks it into standard containers. After that, cranes, trucks and ships need to handle one shape
only, and a new kind of cargo needs a new repacking procedure, not new trucks. The vendable is the
container; the dataset builders are the repacking sheds. The analogy also predicts the vendable's
best-known limitation: a container's label says what is inside, not where it came from on a map,
and the vendable likewise carries no geographic location (\cref{sec:product:georef}).
\end{analogy}

\subsection{Three variants and their fields}

\begin{table}[htbp]\centering
\caption{The three vendable classes.}\label{tab:product:variants}
\small
\begin{tabular}{@{}lp{0.26\linewidth}p{0.25\linewidth}@{}}\toprule
Class & Sensors & Cube holds\\\midrule
\texttt{VendableHyperspectralDataset} & PRISMA, AVIRIS-NG (and SpectralEarth patches) & reflectance\\
\texttt{VendableEnmapHyperspectralDataset} & EnMAP & reflectance, plus five quality masks\\
\texttt{VendableThermalDataset} & Landsat 9, HotSat-1 & Celsius (HotSat-1: raw DN)\\\bottomrule
\end{tabular}
\end{table}

\begin{table}[htbp]\centering
\caption{Fields of \texttt{VendableHyperspectralDataset}
(\texttt{app/models/dataset/vendables.py}, lines 141--178).}\label{tab:product:fields}
\small
\begin{tabular}{@{}llp{0.42\linewidth}@{}}\toprule
Field & Type, shape & Meaning\\\midrule
\texttt{normalized\_hyperspectral\_cube} & array \shape{(C, H, W)} & reflectance; what a model sees\\
\texttt{validity\_cube} & array \shape{(C, H, W)} & 1 real measurement, 0 not\\
\texttt{spectral\_family\_order} & list of $C$ enums & detector of each band (VNIR or SWIR)\\
\texttt{band\_cw\_order} & list of $C$ floats & centre wavelength of each band, nm\\
\texttt{band\_fwhm\_order} & list of floats & FWHM of each band; empty on the common grid\\
\texttt{band\_validity\_by\_position} & list of $C$ ints & per-band validity flag\\
\texttt{band\_level\_validity\_score} & list of $C$ floats & percentage of valid pixels per band\\\bottomrule
\end{tabular}
\end{table}

\begin{figure}[htbp]\centering
\begin{tikzpicture}[font=\footnotesize]
% the two big arrays as stacked slabs
\foreach \k in {4,3,2,1,0}{
  \draw[spec,fill=spec!10] (0.12*\k,0.12*\k) rectangle ++(3.3,2.2);}
\node[align=center,font=\scriptsize] at (1.65,1.1) {\texttt{normalized\_}\\\texttt{hyperspectral\_cube}\\\texttt{(C, H, W)}};
\foreach \k in {4,3,2,1,0}{
  \draw[maskgrey!80!black,fill=maskgrey!12] (4.2+0.12*\k,0.12*\k) rectangle ++(2.6,2.2);}
\node[align=center,font=\scriptsize] at (5.5,1.1) {\texttt{validity\_cube}\\\texttt{(C, H, W)}\\0/1};
% per-band lists
\foreach \y/\t/\c in {2.9/band\_cw\_order/spec,2.25/band\_fwhm\_order/spec,1.6/spectral\_family\_order/spat,0.95/band\_validity\_by\_position/maskgrey!80!black,0.3/band\_level\_validity\_score/maskgrey!80!black}{
  \node[anchor=west,text=\c,font=\scriptsize] at (7.9,\y) {\texttt{\t}};
  \foreach \i in {0,...,7}{\draw[\c] (7.95+0.22*\i,\y-0.4) rectangle ++(0.2,0.16);}
  \node[anchor=west,font=\scriptsize] at (9.75,\y-0.32) {$\times C$};}
\draw[decorate,decoration={brace,amplitude=4pt,mirror}] (0,-0.3) -- (7.3,-0.3)
  node[midway,below=5pt,align=center]{two arrays of the same shape: values and validity};
\draw[decorate,decoration={brace,amplitude=4pt,mirror}] (7.9,-0.3) -- (10.4,-0.3)
  node[midway,below=5pt,align=center]{per-band lists};
\end{tikzpicture}
\caption{The hyperspectral vendable as a data structure: two arrays of identical shape (the
reflectance cube and its validity cube) and five lists with one entry per band. There is no field
for map coordinates.}\label{fig:product:vendable}
\end{figure}

\begin{sourcenote}[: the EnMAP class is not a subclass]
The Indradhanu course (part 4) says EnMAP has ``a subclass'' of
\code{VendableHyperspectralDataset} with five extra quality masks. In the source,
\code{VendableEnmapHyperspectralDataset} is a separate class deriving directly from Pydantic's
\code{BaseModel} (\file{app/models/dataset/vendables.py}, line 233), with the same core fields
re-declared plus the masks, a land/water class layer and the two detectors' channel indices. Code
that checks \code{isinstance(v, VendableHyperspectralDataset)} will therefore \emph{not} accept an
EnMAP onboarding vendable. (After \code{band_filter_apply}, the EnMAP builder's output on the
common grid is the same class as every other hyperspectral vendable.)
\end{sourcenote}

\subsection{Two vendables in a scene's life}

A hyperspectral scene has two vendables. Scene onboarding builds the \textbf{onboarding vendable},
with all the sensor's native bands. The \code{band_filter_apply} Action builds the
\textbf{filtered vendable}, on the common grid, and pickles it to \code{filtered_vendable.pkl}
(\cref{ch:bandpipeline} is that Action's pipeline). Indradhanu v0.2.0 and the classical detectors
read the filtered one; spectral matching reads the onboarding one, because it needs native bands
and their FWHM. If \code{band_filter_apply} has not run, the scoring Action refuses a hyperspectral
scene with \code{"anomaly_scoring on a hyperspectral scene needs input_band_filter_output_id."}
Thermal scenes skip band filtering and are scored from the onboarding vendable.

\coursefig[0.95\linewidth]{iw-04-vendable}{How a vendable reaches Indradhanu (course figure,
Indradhanu part 4). Of the filtered vendable's fields the model uses the cube and the validity
cube; a pixel is spatially valid if it is valid in at least one band.}{fig:product:coursevendable}

\subsection{Per-band and per-pixel validity, computed}

Two summaries of the validity cube $\mat{V}$ appear everywhere. The \emph{per-band} score, stored in
\code{band_level_validity_score}, is the percentage of pixels valid in band $c$:
\begin{equation}\label{eq:product:bandscore}
s_c = 100 \times \frac{1}{HW}\sum_{i=0}^{H-1}\sum_{j=0}^{W-1} \mat{V}[c,i,j].
\end{equation}
The \emph{per-pixel} map used by the scoring code is
\begin{equation}\label{eq:product:spatial}
m(i,j) = \begin{cases} 1 & \text{if } \sum_{c=0}^{C-1}\mat{V}[c,i,j] > 0,\\ 0 & \text{otherwise.}\end{cases}
\end{equation}

\begin{worked}[label={we:product:validity}]{Both summaries on a 3-band, 2-by-2 validity cube}
A validity cube after the band pipeline (all-or-nothing per pixel) has
\[
\mat{V}[0] = \begin{pmatrix}1&1\\0&1\end{pmatrix},\quad
\mat{V}[1] = \begin{pmatrix}1&1\\0&1\end{pmatrix},\quad
\mat{V}[2] = \begin{pmatrix}1&1\\0&1\end{pmatrix}.
\]
\emph{Per band}, by \cref{eq:product:bandscore} with $HW = 4$:
\begin{align*}
s_0 &= 100 \times (1 + 1 + 0 + 1)/4 = 100 \times 3/4 = 75.0,
\end{align*}
and likewise $s_1 = s_2 = 75.0$.

\emph{Per pixel}, by \cref{eq:product:spatial}, sum down the bands at each position:
\begin{align*}
(0,0):\ 1 + 1 + 1 &= 3 > 0 \Rightarrow 1, & (0,1):\ 1 + 1 + 1 &= 3 \Rightarrow 1,\\
(1,0):\ 0 + 0 + 0 &= 0 \Rightarrow 0, & (1,1):\ 1 + 1 + 1 &= 3 \Rightarrow 1.
\end{align*}
So $m = \begin{pmatrix}1&1\\0&1\end{pmatrix}$, and the fraction of valid pixels is $3/4$, the
same $75\%$, as it must be when validity is all-or-nothing.
\end{worked}

\subsection{The contract a model receives}

From the filtered vendable, the scoring code builds exactly two tensors for a hyperspectral model
(\file{_anomaly_scoring_run.py}, lines 285--286):
\begin{lstlisting}
scene_tensor = torch.from_numpy(cube_np).float()                 # (165, H, W) reflectance
mask_tensor  = torch.from_numpy(spatial_valid[None]).float()     # (1, H, W)   1.0 valid
\end{lstlisting}
\code{[None]} inserts a new axis of size $1$ at the front, turning \shape{(H, W)} into
\shape{(1, H, W)}. Notice what is \emph{not} in the contract: no wavelengths, no coordinates, no
sensor name. The model does not know which satellite took the picture; that is the payoff of the
common grid.

\begin{sourcenote}[: a naming collision]
The field is called \code{normalized_hyperspectral_cube}, but ``normalized'' here means ``cleaned
and standardised by the intake pipeline'' (physical units, BSQ layout, common grid). It is
\emph{not} the per-band z-score of \cref{ch:bandpipeline}, which happens inside the model at run
time. The vendable's values are reflectance, around $0.02$ to $0.30$.
\end{sourcenote}

\section{No spatial reference, on purpose}\label{sec:product:georef}

A vendable does not know where on Earth it is. There is no coordinate reference system, no
latitude or longitude, no transform field. To see what is missing, you need one definition.

\begin{definition}[Affine geotransform]\index{affine transform}
A \textbf{geotransform} maps a pixel's column $j$ and row $i$ to map coordinates $(x, y)$ in some
coordinate reference system (CRS):
\begin{equation}\label{eq:product:affine}
\begin{pmatrix} x \\ y \end{pmatrix} =
\begin{pmatrix} a & b \\ d & e \end{pmatrix}
\begin{pmatrix} j \\ i \end{pmatrix} +
\begin{pmatrix} c \\ f \end{pmatrix}.
\end{equation}
$a$ and $e$ are the pixel width and height in map units, $b$ and $d$ are rotation terms (usually
$0$), and $(c, f)$ is the map position of the top-left corner. The \textbf{identity transform} has
$a = e = 1$ and $b = c = d = f = 0$, so $(x, y) = (j, i)$: ``pixel $(0,0)$ is at $(0,0)$ and each
pixel is one unit''. Meaningless as geography.
\end{definition}

The matrix-times-vector product in \cref{eq:product:affine} means: the first output is
$a \cdot j + b \cdot i$, the second is $d \cdot j + e \cdot i$; then add $c$ and $f$.

\begin{worked}[label={we:product:affine}]{Pixel to map, with a real transform and with the identity}
A scene in a UTM projection (map units are metres) has $30$ m pixels, top-left corner at
$x = 500{,}000$, $y = 2{,}000{,}000$, and no rotation (illustrative values\tagA). Rows increase
southwards, so $e = -30$. The transform is $a = 30$, $b = 0$, $c = 500{,}000$, $d = 0$,
$e = -30$, $f = 2{,}000{,}000$. Where is the pixel at row $i = 10$, column $j = 20$?
\begin{align*}
x &= a\,j + b\,i + c = 30 \times 20 + 0 \times 10 + 500{,}000 = 600 + 0 + 500{,}000 = 500{,}600,\\
y &= d\,j + e\,i + f = 0 \times 20 + (-30) \times 10 + 2{,}000{,}000 = -300 + 2{,}000{,}000 = 1{,}999{,}700.
\end{align*}
With the identity transform the same pixel lands at $(x, y) = (20, 10)$: twenty metres east and ten
north of wherever $(0,0)$ is in the file's CRS. Opened in a GIS, an intermediate GeoTIFF therefore
lands as a tiny image near the origin (for a geographic CRS, off the coast of West Africa). That is
not a bug.
\end{worked}

Every GeoTIFF an Action writes carries the identity transform. The real georeferencing is
recovered once, at export, by re-reading the original raw file (\code{app/georef/}). Carrying
geographic metadata through every step would mean every step had a chance to corrupt it; deriving it
once, from the source, at the end, means it is right. PRISMA is the awkward case: its L2D product is
swath data with a latitude and longitude per pixel, and the export fits an axis-aligned
affine transform to those, discarding the swath's curvature (\file{docs/03-data-pipeline.md}). A
scene whose georeferencing cannot be resolved fails at export, not at analysis, with the error
\code{422 crs_missing}.

Practical consequences: do not add a transform field to the vendable expecting it to be filled in,
and do not trust the coordinates of an intermediate \code{.tif}.

\section{Vendables are pickled: class names are a wire format}\label{sec:product:pickle}

\code{band_filter_apply} writes the filtered vendable to disk with Python's \code{pickle} module,
and later Actions (and the api) read it back. A pickle stores an object's data together with the
\emph{full import path} of its class, for example
\code{app.models.dataset.vendables.VendableHyperspectralDataset}. Loading the file imports that
exact path. Rename the class or move the module, and every vendable already on disk becomes
unloadable: not corrupted, unreadable, because the recorded path no longer exists. The module's own
docstring says it: ``Treat the class names and module path as a wire format.'' Adding a field is
safe; renaming or moving one is a breaking change, like changing a database schema.

\section{The thermal vendable and its two cloud masks}\label{sec:product:thermal}

\code{VendableThermalDataset} holds \code{normalized_thermal_cube} (degrees Celsius) and a
validity cube, plus two kinds of cloud information kept apart on purpose: a \code{cloud_mask}
modelled from the scene's own temperature distribution (a five-component Gaussian mixture,
\file{docs/03-data-pipeline.md}) and the provider's \code{QA_PIXEL}-derived masks
(\cref{we:sensors:qa}). The provider's mask is authoritative but conservative and not always
present; the modelled one is always available and adapts to the scene. Merging them would throw
away their disagreement, which is often the interesting part.

\begin{sourcenote}[: HotSat-1 units are silently dropped]
The anomaly-detection course (part 5) notes, and the source confirms, that the HotSat-1 builder
passes a \code{units} value that \code{VendableThermalDataset} does not declare. Pydantic models
ignore undeclared fields by default, so the value is discarded, and the scoring code's check
\code{getattr(vendable, "units", None)} finds nothing. The guard meant to treat HotSat's raw DN
differently from Celsius (\cref{ch:bandpipeline}) is therefore not in effect through this path.
Source over comments.
\end{sourcenote}

\begin{confusions}
\confusion{Can one Project cover several Scenes?}{No. A Project binds to exactly one Scene. Two
scenes, two Projects.}
\confusion{Is the Action graph a workflow engine?}{Lighter than that. Actions name their inputs;
there is no scheduler and no automatic re-run. A user creates each Action.}
\confusion{Why is scoring separate from thresholding?}{Scoring is deterministic and expensive;
thresholding is a judgement call and cheap. Splitting them lets you try several thresholds on one
scoring run.}
\confusion{Is Indradhanu a service I can call?}{No. It is a Python class loaded inside the worker,
on demand, for the duration of one scoring job, as one of the models ticked in
\texttt{anomaly\_scoring}.}
\confusion{Is the vendable the file on disk or the object in memory?}{Both, at different moments.
It is a Python object; \texttt{filtered\_vendable.pkl} is that object pickled.}
\confusion{The intermediate GeoTIFF sits at (0,0). Is georeferencing broken?}{No. Action outputs
carry the identity transform by design; coordinates are recovered at export from the raw file.}
\confusion{The api is healthy but nothing runs. Why?}{Almost always the worker: crashed or
crash-looping, often on an import error. The api's health check does not look at the worker.}
\end{confusions}

\begin{checkyourself}
\item\label{ex:product:1} A user wants to analyse one scene with two band-filter settings, score each
with Indradhanu and MNF-RX in one scoring Action, and try two thresholds on each scoring run. How many
Scenes, Projects and Actions (counting only \texttt{band\_filter\_apply}, \texttt{anomaly\_scoring}
and \texttt{anomaly\_detection\_prep}) are needed if each band-filter setting has its own Project?
\item\label{ex:product:2} For the graph with edges $B \to G$, $B \to S$, $G \to S$, $S \to P$,
$P \to M$, $S \to Q$, write the in-degree of each node and give two different topological orders.
Explain why adding an edge $M \to B$ is impossible in the product.
\item\label{ex:product:3} A validity cube of shape $(2, 2, 3)$ has $\mat{V}[0] = \mat{V}[1] =
\begin{pmatrix}1&0&1\\1&1&0\end{pmatrix}$. Compute $s_0$ and the per-pixel map $m$, and the fraction
of valid pixels.
\item\label{ex:product:4} With the transform of \cref{we:product:affine}, where is pixel row $i = 0$,
column $j = 0$? Row $100$, column $0$? And which pixel is at map position
$x = 501{,}500$, $y = 1{,}998{,}500$? (Invert \cref{eq:product:affine} by hand.)
\item\label{ex:product:5} With heartbeat interval $5$ s and timeout $60$ s, a worker freezes for
$45$ seconds and then resumes. Is its job reaped? What if the operator sets the timeout to $12$ s?
\end{checkyourself}

\nextchapter{ch:bandpipeline}
